
\chapter{Conclusão}

A presente prática laboratorial cumpriu integralmente seus propósitos didáticos e técnicos ao transitar pela modelagem algébrica simbólica e culminar na simulação de circuitos de potência chaveados. A interpretação dos resultados numéricos revela que \textbf{os requisitos de projeto do conversor \textit{Buck} foram plenamente atendidos}. Em todas as instâncias simuladas, a ondulação $\Delta V$ permaneceu inferior ao teto de $5,0\text{ mV}$, comprovando que as expressões deduzidas via expansão de Taylor em alta frequência ($T \to 0$) fornecem um dimensionamento de indutores altamente seguro e eficaz.

Analisando a exatidão das tensões de saída ($V_0$), observa-se um comportamento perfeitamente alinhado à teoria nos Cenários 2, 3 e 4, com erros percentuais irrisórios (abaixo de $0,2\%$). No entanto, o Cenário 1 ($R = 5\,\Omega$, $L = 125\text{ mH}$) apresentou um desvio fora do esperado, estabilizando a tensão em cerca de $2,94\text{ V}$ ao invés dos $2,50\text{ V}$ nominais.

Essa anomalia aparente pode ser tecnicamente justificada pela constante de tempo do circuito ($\tau = L/R$). No Cenário 1, $\tau = 125\text{ mH} / 5\,\Omega = 25\text{ ms}$. Sabendo que um sistema de primeira ordem requer aproximadamente $4\tau$ a $5\tau$ (cerca de $100\text{ a }125\text{ ms}$) para atingir integralmente o regime permanente, a janela de simulação transiente encerrada em $100\text{ ms}$ capturou o circuito ainda na cauda do seu amortecimento final, antes da completa acomodação térmica e elétrica. Adicionalmente, grandes indutores operando com baixas correntes de carga tendem a ressaltar as não-idealidades mínimas de condução e corte do diodo de roda livre modelado no SPICE.

Por fim, o experimento validou a metodologia de aproximação linear por partes para o trato de componentes não-lineares, atestando que a Transformada de Laplace, aliada ao suporte algébrico de softwares como o Maxima \cite{maxima}, configura-se como uma ferramenta poderosa e insubstituível para o domínio de transitórios em engenharia eletrônica.